\color{unidarkgreen}{\section[(рекомедательное) Схема подключения ИЭУ для случаев подключения АТ]{(рекомендательное)\\Схема подключения ИЭУ для случаев подключения АТ}\label{app:schema}}
\color{black}

\begin{figure}[H]
\centering

% Первая строка: а) и б)
\begin{subfigure}{0.45\textwidth}
    \centering
    \includegraphics[width=\textwidth, trim=0mm 0mm 130mm 0, clip]{\apppath/Приложение. Схема подключения/_latex/img/Fig 1.pdf}
    \caption{а)}
    \label{pict:appschema_1}
\end{subfigure}
% Пробел между а) и б) — 5% ширины
\hfill
\begin{subfigure}{0.43\textwidth}
    \centering
    \includegraphics[width=\textwidth, trim=135mm 0mm 0 0, clip]{\apppath/Приложение. Схема подключения/_latex/img/Fig 1.pdf}
    \caption{б)}
    \label{pict:appschema_2}
\end{subfigure}

% Добавляем вертикальный отступ с сохранением центрирования
\par\vspace{4pt}

% Вторая строка: в) — центрируем её относительно общей ширины
\begin{subfigure}{0.4\textwidth}
    \centering
    \includegraphics[width=\textwidth, trim=0mm 0mm 0 0mm, clip]{\apppath/Приложение. Схема подключения/_latex/img/Fig 2.pdf}
    \caption{в)}
    \label{pict:appschema_3}
\end{subfigure}

% Общая подпись с выравниванием пунктов а, б, в
\captionsetup{justification=raggedright, singlelinecheck=off}
\caption{%
  Поясняющие схемы подключения шкафа для случаев подключения автотрансформатора: \newline
  \hspace*{1.5em}а) по схеме одна рабочая секционированная система шин с подключением трансформаторов через развилку из выключателей 220-9Н (комплект КСЗ установлен на стороне ВН АТ);  \newline
  \hspace*{1.5em}б) по схеме две рабочие и обходная система шин (комплект КСЗ установлен на стороне СН АТ); \newline
  \hspace*{1.5em}в) по схеме трансформатор-шины с полуторным присоединением линии (комплект КСЗ установлен на стороне ВН АТ).
}
\label{pict:schema}
\end{figure}

\medskip
\clearpage